\section{Discussion and Threats to Validity}
\label{sec:discussion}

\textbf{Reading ``guessed'' correctly.} SemLoc's top line is the LLM's guess by design, and the failure analysis shows the
guess is often semantically right and one line off the label. The lesson is therefore not that the bonus should go (on
our runs removing it changes Acc@1 by \RqTwoRepairedAllSemlocNoPriorDAccOne{} points) but that a tool which claims to be
grounded should state that its first answer is the model's, and that line-level labels need to accept equivalent fix
locations before small differences between methods mean much.

\textbf{Why tests make poor arbiters of location.} Three results share one cause. SemLoc's counterfactual step changes no
ranking; our minimal-fix oracle moved \ReviewCorrectedOverturnedN{} of its \GtCorrectedN{} relocated labels to lines the review
rejected; and propose-and-verify does not improve on its proposer. All three accept a line when a small patch there makes
every test pass. With about ten tests per program, that criterion holds in \PavmechReviewedOtherLineFixable\% of
fix attempts at non-faulty candidate lines, usually through a compensating constant or subtraction, and in only
\PavmechReviewedFaultLineFixable\% at faulty ones (dev split, the proposer's top candidates), which two-line faults, omissions and larger edits put out of reach. This is the plausible-patch
problem~\cite{qi2015kali,smith2015cure} seen from localization: the tests suffice to \emph{suggest} where the bug is, not to
\emph{prove} it.

\textbf{Saturation and practice.} With the reviewed (LLM-suggested) labels one prompt reaches \RqTwoReviewedAllLlmTestsAccThreeMean\% Acc@3.
SemFault-250's programs are single short functions, its faults are small local mutations, and its test names state the
violated intent (e.g., \texttt{test\_fee\_deducted\_once\_not\_per\_participant}); much of the difficulty the original paper
measured came from mislabels. The benchmark can no longer show that an elaborate pipeline is needed; real bugs, where
every method stays far from the ceiling, can. For practitioners, the strongest investment is tests that state intent, as
test-driven development produces; and an AI agent that stops when the tests pass should have its patch reviewed, since
about a quarter of fix attempts at plausible but non-faulty lines yield a test-passing patch on suites of this size.

Table~\ref{tab:checklist} collects seven checks for evaluating LLM-based fault localization; each changed at least one number in
this study.

\begin{table}[!t]
\centering
\caption{Checklist for evaluating LLM-based fault localization, with what each check changed here}
\label{tab:checklist}
\scriptsize
\begin{tabularx}{\columnwidth}{@{}>{\raggedright\arraybackslash}p{0.38\columnwidth}>{\raggedright\arraybackslash}X@{}}
\toprule
Check & Effect in this study \\
\midrule
Average-rank ties & SBFL Acc@5 \RqOneSbflOchiaiAuthorAllAccFive\% $\to$ \RqOneSbflOchiaiFairAllAccFive\% \\
Fixed denominator, any-hit not rank one & delta debugging \RqOneDdAuthorAllAccOne\% $\to$ \RqOneDdFairAllAccOne\% in the released script \\
Checked labels & \LabaudAllMoved{} of \LabaudAllN{} labels move; SemLoc's Acc@3 lead over LLM-Direct vanishes on reviewed labels \\
Ablate every component, incl.\ undocumented & guess bonus: Acc@1 changes by \RqTwoRepairedAllSemlocNoPriorDAccOne{} without it \\
Repeats with paired CIs & several single-repeat differences are not significant; pooled ones are \\
Report cost & counterfactual step: \CostCfSharePct\% of calls, no ranking change \\
Single prompt with failing tests & beats the pipeline at Acc@1 on all label sets (\RqTwoRepairedAllLlmTestsDAccOne{} repaired) \\
\bottomrule
\end{tabularx}
\end{table}

\textbf{Threats to validity.}\label{sec:threats} Internally, the minimal-fix oracle can accept compensating fixes and never searches
alternatives for reachable labels; the review overturned \ReviewCorrectedOverturnedN{} of its relocations. The reviewed
labels are LLM-suggested; in the first round reviewers saw the suggestion and agreed in \ReviewAgreeN{} of \ReviewHumanN{}
cases. The blind re-review supports them, but on a sample of \BlindBlindAllPrograms{} programs, with
\BlindAnswersExcluded{} non-blind answers excluded and choices limited to the listed candidates, so the reviewed labels are
still not fully independent of LLM judgment and may favor LLM-based methods; we therefore state main results on the original and
repaired labels, which involve no LLM judgment of location. Two-line labels (\LabaudAllTwoLine{} programs) make the reviewed
Acc@$k$ slightly easier. Line-level Acc@1 depends on whether a guard or the guarded statement is labeled; we report tolerant
variants. The single prompt is ours, fixed before seeing test results. The knobs of propose-and-verify were selected on the
dev split and then applied unchanged to the test split.

Externally, test names in SemFault-250 state intent unusually clearly; on suites with generic names the
single prompt's advantage may shrink, and we did not run a hidden-test-name ablation. Every program is one function, and
for real bugs the target function is taken from the developer fix, so localizing the function within a project, the
hard part on large code bases, is not evaluated. Real bugs come from one project (31 youtube-dl bugs). RQ2, RQ3, the
decomposition and the renaming study use one model; the model-family subset has no CIs, and GPT-6 luna's FailLoc runs used the OpenAI API rather than Codex; models run at provider-default
reasoning. Renaming changes identifiers and comments but not structure or test names, so memorization of templates is
not fully excluded.
